\section{密結合方式による軌道計算}
\label{sec:coupling}

\subsection{方式の選定}

検討したいピッチ条件は少数（数条件程度）であるため，軌道再現の忠実度を優先し，係数テーブルによる補間方式（疎結合方式）ではなく，\textbf{CFDと運動方程式を毎ステップ結合する密結合方式}を採用する。少数条件のみを高精度に再現したい場合，密結合方式のコストは疎結合方式の2--3倍程度に留まる一方，条件数が数十を超える多条件スタディでは疎結合方式が圧倒的に有利になる（\(10^4\)条件規模で約200倍のコスト差）。将来的に多条件のパラメトリックスタディが必要になった場合は，疎結合方式への切り替えを推奨する。

\subsection{ボール追従座標系（非慣性系）}
\label{sec:frame}

密結合とは，マウンドからホームベースまでの約18.4\,m（ボール約400個分）を実空間で移動格子等により直接解くことを意味しない——単一のコンシューマGPU（8\,GB VRAM）でこれを行うのは非現実的である。そこで\textbf{ボール追従非慣性座標系（body-fitted moving reference frame）}を採用し，計算領域自体はボール周囲数直径角の小さな箱に保ったまま，実時間で運動方程式と結合する。

\begin{enumerate}
	\item 計算領域はボールの重心を原点とする座標系に固定する（ボールは領域内で自転するのみで並進しない）。
	\item 流入境界条件は，その瞬間のボールの対気速度ベクトル \(-\vec V_{ball}(t)\) に相当する一様流として毎ステップ更新する。
	\item 座標系が加速度運動（重力による落下，抗力による減速）をしているため，非慣性系補正として一様体積力（見かけの重力項）をLBMのフォーシング項に追加し，流体側の運動量収支を物理的に正しく保つ。
	\item 各タイムステップ（またはサブサイクルごと）で運動量交換法（\secref{geometry-boundary}）により瞬時の \(\vec F_{aero}\)，\(\vec T_{aero}\) を算出し，6自由度運動方程式（\eref{translation}, \eref{rotation}）に代入してその場でRK4により時間積分する。
	\item 更新された対気速度・角速度を，次ステップの流入境界条件・体積力・自転境界条件へフィードバックする。
\end{enumerate}

この構成により，係数テーブルを介さずCFDと運動方程式が毎ステップ双方向に結合された「真の密結合」でありながら，計算領域を小さな箱に保てるため，単一の8\,GB級GPUでも実行可能になる。

\subsection{見かけの体積力の導出}
\label{sec:apparent-force}

フレームはボールと同じ加速度
\begin{align}
	\vec A = \frac{d\vec V}{dt} = \frac{\vec F_{aero}}{m} + \vec g
	\label{eq:frame-accel}
\end{align}
で運動するため，流体の運動方程式には一様な見かけの力 \(-\rho\vec A\) が加わる。一方，実際の重力 \(+\rho\vec g\) は，静止大気中では静水圧勾配と釣り合っているため，両方とも落とすことと両方とも持つことが等価である（差は浮力のみで，野球ボールでは2.6\,mN対自重1.42\,N＝0.18\%）。したがって流体が受ける一様体積加速度は
\begin{align}
	\vec a_{fluid} = -\vec A = -\left(\frac{\vec F_{aero}}{m} + \vec g\right)
	\label{eq:apparent-force}
\end{align}
となり，\textbf{重力は流体側には直接現れない}（ボールの応答を通してのみ流れ場に効く）。自由落下（\(\vec F_{aero}=0\)）では \(\vec a_{fluid}=-\vec g\)，すなわち空気が上向きに \(g\) で加速して見える，という直感とも一致する。

この形が本質的に重要なのは，次の恒等式が成り立つためである。フレームから見た主流速度は \(\vec u_\infty=-\vec V(t)\) であり，
\begin{align}
	\frac{d\vec u_\infty}{dt} = -\vec A = \vec a_{fluid}.
	\label{eq:identity}
\end{align}
つまり\textbf{フレームの加速度を表す体積力そのものが，遠方場を次ステップの流入条件が要求する速度へちょうど運ぶ}。境界条件と内部場が食い違って静かに破綻する，という移動座標系スキームの典型的な失敗モードが構造的に起きない。

6自由度側（並進・自転・姿勢）では，姿勢はジンバルロック（スピン軸が飛行方向に近づくジャイロ成分の強い変化球で現実的な危険）を避けるため四元数で保持し，積分にはRK4を用いる。CFDから凍結した力を与える限り，並進・自転についてRK4は厳密である。

\subsection{密結合ループとサブサイクル平均}
\label{sec:coupling-loop}

密結合ループの1回の外側反復は次の4段からなる：

\begin{enumerate}
	\item 箱平均速度を測り体積力（\eref{apparent-force}）を決める。
	\item サブサイクル分だけLBMを進め，運動量交換の力をその区間で平均する。
	\item 力・トルクをN・N\(\cdot\)mに変換し，6自由度運動方程式を1ステップ進める。
	\item 更新された速度・角速度を(1)へ戻す。
\end{enumerate}

\paragraph{サブサイクルの目的は省力化ではなく平均化である}
ボールの速度は1球（約\(1.85\times10^5\)ステップ）で15\%しか変わらない——1ステップあたり\(10^{-6}\)——であり，軌道側は毎ステップの更新を必要としない。一方，乱流中の物体に働く瞬時の運動量交換力は平均値より大きく揺らぐため，区間平均でなければ軌道に与える標本として不適切である。

\paragraph{力の使い方に一貫性が要る}
同一ステップ内では，直前のサイクルで測った力を体積力にも6自由度積分にも使う。\eref{identity}の恒等式が項ごとに成立するのは，フレーム加速度を表す力とボールを加速する力が\emph{同じ数値}であるときだけである。新しく測った力を軌道に，前サイクルの力を流体に使うと，遠方場と流入条件が1サイクル分ずれる。

\subsection{力の分解}

各サブサイクルで運動量交換の力を \((\hat u,\hat n,\hat b)\) 系に分解して記録する：\(\hat u=\vec V/|\vec V|\)（抗力方向），\(\hat n=(\vec\omega\times\hat u)/|\vec\omega\times\hat u|\)（マグヌス方向。飛行方向成分のスピン＝ジャイロ成分は構成上自動的に落ちる），\(\hat b=\hat u\times\hat n\)。第3成分 \(C_{side}\) が本プロジェクトの要点である（\secref{magnus}）：実球は縫い目がマグヌス面まわりの圧力分布を非対称にするためこの方向に力を持ち，それがスウィーパーやジャイロボールが教科書どおりに曲がらない主因である。解析的な係数モデルはここに必ずゼロを返すが，CFDは測った値を返せる。

\subsection{流入・流出境界条件}
\label{sec:open-boundary}

計算領域を単純な周期境界とした場合，ボールは自分自身の後流の中を飛ぶことになる。抗力係数の基準とすべき「球に届く流れ」が周期箱では曖昧になり，その曖昧さが本来測ろうとしている物理量の差より大きくなりうることが検証（\secref{validation}）で判明したため，流れ方向の面だけを非周期の流入・流出境界に置き換える機能を実装した（\texttt{src/boundary/open.jl}）。

流れ方向は \(-x\) 方向であるため，高\(x\)面が入口，低\(x\)面が出口となる。いずれもGuoの非平衡外挿を用い，内側の第1平面から非平衡部 \(f^{neq}\) を持ってくる：

\begin{itemize}
	\item \textbf{入口}：速度を課し，密度は内側から取る。自由流を「推定」せず「指定」する。
	\item \textbf{出口}：密度を課し，速度は内側から取る。後流が対流で流出し，圧力の基準がここで固定される。
\end{itemize}

分担が逆だと成立しない。出口が圧力を固定しないと箱全体の密度が漂い，入口が密度まで課すと過剰決定になって条件が跳ね返る。

流入・流出境界を導入すると，周期境界で必要だった平均流コントローラ（\secref{controller}）は不要になる——入口があれば自由流を直接指定でき，ボール追従座標系では\eref{identity}の恒等式が「内部の流体は入口値と同じ率で加速する」ことを保証するので，フィードバックそのものが要らない。

\subsection{平均流コントローラ（周期境界での暫定措置）}
\label{sec:controller}

流入・流出境界を用いない構成（完全周期境界）では，球が受ける抗力の分だけ流体の運動量が行き場を失い，主流が減衰する。そこで箱平均速度を目標値 \(-\vec V(t)\) に保つ一様体積力を加える（constant-mass-flux forcing）。

このコントローラはPI制御でなければならない。純積分制御では誤差が \(\ddot e=-Ke\) に従い減衰しないため体積力が目標値のまわりで振動するが，目標値自体が極めて小さい（定常補正は1サブサイクルあたり平均速度を0.03\%程度動かす量）ため，わずかな振動でも信号が埋もれてしまう。比例項を加えて \(K_p=2/T\)，\(K_i=1/T^2\)（臨界減衰）とすると運動量収支の残差は実用的な水準に収まる。ここで \(T\) は箱の通過時間の数倍にとる。

定常状態では，コントローラが供給する運動量と表面から運動量交換法で得た力が一致しなければならない（フレームの体積力は\eref{identity}の恒等式により相殺して現れない）。この2つはまったく別経路で計算されるため，両者の残差は強制項・境界条件・単位換算のすべてが整合していることの常時チェックとして使える。

この構成には既知の限界がある：周期像によるブロッケージは残るため，領域幅を十分にとる必要がある。検証（\secref{validation}）により，この限界は想定より重いことが定量的に判明している。
